\documentclass[11pt,a4paper]{article}

%% =========================================================================
%% PACKAGES & PREAMBLE
%% =========================================================================
\usepackage[utf8]{inputenc}
\usepackage[T1]{fontenc}
\usepackage[margin=0.85in, top=0.9in, bottom=0.9in, headheight=14pt]{geometry}
\usepackage{amsmath,amssymb,amsfonts,mathtools}
\usepackage{booktabs,tabularx,multirow,array}
\usepackage{graphicx,xcolor}
\usepackage{hyperref}
\usepackage{fancyhdr}
\usepackage{enumitem}
\usepackage{listings}
\usepackage{tikz}
\usetikzlibrary{shapes.geometric, arrows.meta, positioning, calc, backgrounds}
\usepackage{caption}
\usepackage{subcaption}
\usepackage{titlesec}
\usepackage{cite}
\usepackage{microtype}
\usepackage{float}
\usepackage[ruled,vlined]{algorithm2e}

%% Hyperref configuration
\hypersetup{
    colorlinks=true,
    linkcolor=teal!80!black,
    citecolor=blue!80!black,
    urlcolor=teal!70!black,
    pdfauthor={Maloth Madhu},
    pdftitle={AI-based Village Pond Planning System - Final Technical Report},
    pdfkeywords={Village Pond Planning, Geospatial Analysis, D8 Flow Routing, SCS-CN Model, Web GIS, FastAPI, IIT Bhilai}
}

%% Header & Footer configuration
\pagestyle{fancy}
\fancyhf{}
\fancyhead[L]{\small\textbf{CSD Assignment 1 -- Phase 2 Report}}
\fancyhead[R]{\small\textbf{JalDrishti: Village Pond Planning System}}
\fancyfoot[L]{\small\textit{Maloth Madhu (ID: 12341370) -- IIT Bhilai}}
\fancyfoot[R]{\small Page \thepage}
\renewcommand{\headrulewidth}{0.5pt}
\renewcommand{\footrulewidth}{0.4pt}

%% Section styling
\titleformat{\section}{\large\bfseries\color{teal!80!black}}{\thesection.}{0.5em}{}[\titlerule]
\titleformat{\subsection}{\normalsize\bfseries\color{teal!70!black}}{\thesubsection}{0.5em}{}
\titleformat{\subsubsection}{\small\bfseries\color{black}}{\thesubsubsection}{0.5em}{}

%% Listings configuration for code blocks
\definecolor{codegray}{rgb}{0.96,0.96,0.96}
\definecolor{codegreen}{rgb}{0.13,0.54,0.13}
\definecolor{codeblue}{rgb}{0.10,0.20,0.70}
\definecolor{codepurple}{rgb}{0.50,0.00,0.50}

\lstset{
    backgroundcolor=\color{codegray},
    basicstyle=\ttfamily\scriptsize,
    breaklines=true,
    captionpos=b,
    commentstyle=\color{codegreen},
    keywordstyle=\color{codeblue}\bfseries,
    stringstyle=\color{codepurple},
    numbers=left,
    numberstyle=\tiny\color{gray},
    numbersep=5pt,
    frame=single,
    rulecolor=\color{black!15},
    tabsize=2,
    showstringspaces=false
}

%% Set paragraph spacing
\setlength{\parskip}{4pt}
\setlength{\parindent}{0pt}

%% =========================================================================
%% DOCUMENT METADATA
%% =========================================================================
\title{
    \vspace{-1.2cm}
    \textbf{\Large INDIAN INSTITUTE OF TECHNOLOGY BHILAI}\\
    \textbf{\normalsize DEPARTMENT OF COMPUTER SCIENCE AND ENGINEERING}\\[0.3cm]
    \rule{\linewidth}{1.2pt}\\[0.3cm]
    \textbf{\LARGE JalDrishti: An AI-Assisted Geospatial Decision Support System for Automated Village Pond Site Selection, Catchment Delineation, and Runoff Estimation}\\[0.2cm]
    \large\textbf{Computer System Design (CSD) -- Assignment 1: Phase 2 Final Technical Report}\\[-0.2cm]
    \rule{\linewidth}{1.2pt}
}

\author{
    \textbf{Maloth Madhu}\\
    \textbf{Student ID No.:} 12341370 \quad $\vert$ \quad \textbf{Email:} \texttt{malothm@iitbhilai.ac.in}\\
    \textbf{Allocated System ID:} \texttt{stu29\_sys2} \quad $\vert$ \quad \textbf{Server Node:} \texttt{10.1.75.51:2314}\\
    \textbf{Department of Computer Science and Engineering}\\
    \textbf{Indian Institute of Technology Bhilai, Durg, Chhattisgarh 491001, India}
}
\date{\textbf{Submission Date:} September 30, 2026}

%% =========================================================================
%% BEGIN DOCUMENT
%% =========================================================================
\begin{document}

\maketitle
\thispagestyle{fancy}

%% =========================================================================
%% ABSTRACT & KEYWORDS
%% =========================================================================
\begin{abstract}
\noindent
Persistent water scarcity during the post-monsoon dry season severely impedes rural agricultural sustainability, livestock sustenance, and domestic water security across semi-arid regions of India. While decentralized rainwater harvesting through excavated village ponds represents a cost-effective, proven intervention aligned with national priorities such as \textit{Mission Amrit Sarovar} and the \textit{Jal Jeevan Mission}, traditional pond site selection remains overwhelmingly reliant on manual, heuristic, and unscientific practices. Such practices frequently overlook fine-grained topographic elevation gradients, micro-watershed boundaries, and multi-year precipitation variability, leading to widespread pond failures characterized by inadequate catchment inflow or rapid desiccation.

To solve this critical challenge, this report presents the engineering design, algorithmic formulation, and production cluster deployment of \textbf{JalDrishti}, an end-to-end, high-performance geospatial decision support platform. JalDrishti enables village panchayat engineers, agricultural extension workers, and hydrologists to draw custom village land boundaries directly onto an interactive Web GIS satellite interface or upload standard vector contour maps (KML/KMZ). Within seconds, the system executes an automated spatial computing pipeline: (1) extracting 3D vector coordinates to reconstruct a continuous Digital Elevation Model (DEM) via bivariate spline interpolation, (2) delineating the exact upstream contributing catchment boundary using a Deterministic Eight-Node (D8) steepest-descent flow direction and accumulation algorithm, (3) estimating annual harvestable surface runoff using the USDA Soil Conservation Service Curve Number (SCS-CN) hydrological framework, and (4) pinpointing the mathematically optimal pond location along with engineering dimension specifications (surface area, target depth, and active volumetric capacity).

The system is architected around a strict separation of concerns, featuring a 100\% vanilla HTML5/CSS3/JavaScript single-page frontend (with zero React and zero WebSockets) connected to an asynchronous Python 3.12 FastAPI backend through an embedded multi-threaded reverse proxy. Deployed permanently on allocated cluster system \texttt{stu29\_sys2} (\texttt{10.1.75.51:2314}), the platform achieves sub-220\,ms end-to-end analysis latency, 24/7 autonomous watchdog self-healing, and complete immunity to campus router packet dropping. Comprehensive experimental evaluations across Chhattisgarh agrarian watersheds demonstrate that JalDrishti can accurately identify natural hydrological convergence zones and calculate harvestable runoff yields exceeding 64\,million liters per square kilometer, thereby providing a robust, democratized platform for rural climate resilience.
\end{abstract}

\textbf{Keywords:} Village Pond Planning, Watershed Delineation, D8 Flow Routing, Digital Elevation Model (DEM), SCS-CN Runoff Model, Web GIS, FastAPI, Reverse Proxy, Asynchronous Systems, Fault-Tolerant Daemons, IIT Bhilai Cluster.

\vspace{0.2cm}
\hrule
\vspace{0.2cm}

%% =========================================================================
%% SECTION 1: INTRODUCTION & PROBLEM STATEMENT
%% =========================================================================
\section{Introduction and Problem Context}
\label{sec:intro}

\subsection{Background and Real-World Motivation}
India supports approximately 18\% of the world's human population and 15\% of global livestock on merely 4\% of the planet's renewable freshwater resources. Over 68\% of the nation's net sown agricultural area is vulnerable to drought, with smallholder and marginal farmers suffering severe crop failures during prolonged dry spells between monsoon events. In agrarian states such as Chhattisgarh, Madhya Pradesh, and Maharashtra, decentralized water harvesting structures---colloquially termed \textit{village ponds}, \textit{talabs}, or \textit{check dams}---form the hydrological backbone of the rural economy. When engineered correctly, village ponds fulfill three critical ecological and economic functions:
\begin{enumerate}[noitemsep,topsep=2pt]
    \item \textbf{Groundwater Replenishment:} Sustained pond water head creates a localized recharge mound, elevating water tables in nearby shallow open wells and tube wells.
    \item \textbf{Supplemental Protective Irrigation:} Stored surface runoff provides critical lifeline irrigation during the post-monsoon \textit{Rabi} cropping cycle, dramatically improving crop yields.
    \item \textbf{Livestock and Domestic Water Security:} Ponds supply non-potable domestic washing, livestock hydration, and freshwater pisciculture opportunities for rural livelihoods.
\end{enumerate}

Recognizing this potential, the Government of India launched the ambitious \textbf{Mission Amrit Sarovar} in 2022, targeting the construction and rejuvenation of at least 75 water bodies in every district of the country. However, the execution of this mission has exposed a fundamental technological gap: village-level administrative bodies (\textit{Gram Panchayats}) and rural engineering staff lack access to computerized, scientifically rigorous hydrological decision support tools.

\subsection{Critical Bottlenecks in Conventional Manual Pond Siting}
Traditionally, village pond construction sites are selected based on intuitive visual inspections, historical precedent, or political boundary conveniences rather than rigorous terrain analysis. This conventional approach suffers from four chronic failure modes:
\begin{enumerate}[noitemsep,topsep=2pt]
    \item \textbf{Inadequate Catchment Drainage Inflow:} Constructing an excavation outside the natural watershed drainage channel results in inadequate overland flow accumulation, leaving the pond partially dry even following heavy monsoon deluges.
    \item \textbf{Premature Siltation and Soil Permeability Pitfalls:} Siting ponds on excessively permeable sandy strata or upstream of severe erosive gullies leads to either uninhibited subterranean infiltration losses or rapid reservoir siltation within 2--3 seasons.
    \item \textbf{Prohibitive GIS Software Complexity:} Commercial Geographic Information System (GIS) suites such as ESRI ArcGIS, Bentley WaterCAD, or even open-source desktop software such as QGIS and GRASS GIS impose steep learning curves, require high-end local workstations, and demand specialized domain training unavailable to village-level field officers.
    \item \textbf{Static, Disconnected Meteorological Data:} Historical precipitation records, evaporation metrics, and land-use rasters are rarely synthesized in real time, preventing dynamic scenario evaluations based on varying monsoon intensities.
\end{enumerate}

\subsection{Project Scope and Deliverables}
The primary objective of the \textbf{JalDrishti} project is to eliminate these barriers by developing and deploying a lightweight, web-accessible, automated decision support system. The scope of this project includes:
\begin{itemize}[noitemsep,topsep=2pt]
    \item Providing a zero-install, responsive Web GIS frontend supporting interactive land polygon drawing over high-resolution satellite imagery (Esri World Imagery, OpenStreetMap, OpenTopoMap, and CartoDB Dark).
    \item Ingesting both interactive map coordinates and external 3D contour vector files (KML/KMZ).
    \item Automating terrain surface modeling, D8 flow direction computation, and exact upstream watershed boundary delineation.
    \item Calculating harvestable annual surface runoff volume using the empirical USDA Soil Conservation Service Curve Number (SCS-CN) method.
    \item Outputting precise engineering recommendations: optimal pond geographic coordinates, excavation depth, surface footprint dimensions, and household water security impact metrics.
    \item Implementing a high-availability, fault-tolerant production deployment on the designated IIT Bhilai cluster node \texttt{stu29\_sys2} (\texttt{10.1.75.51:2314}) accessible through designated port-forwarding mappings.
\end{itemize}

\textit{Explicit Non-Goals:} The platform strictly focuses on hydrological site identification and preliminary reservoir sizing; detailed structural geotechnical engineering (e.g., concrete reinforcement schedules, sub-surface bore-hole permeability logging, and masonry spillway hydraulic design) remains outside the scope of this software system.

%% =========================================================================
%% SECTION 2: THEORETICAL FOUNDATIONS & RELATED WORK
%% =========================================================================
\section{Theoretical Foundations and Related Work}
\label{sec:theory}

\subsection{Terrain Surface Modeling and Digital Elevation Models (DEM)}
A Digital Elevation Model (DEM) is a regular raster grid representing bare-earth topographic elevations $z = f(x, y)$ referenced to a vertical geodetic datum. Topographic vector contour maps, typically distributed as Keyhole Markup Language (KML/KMZ) files, represent elevations as discrete iso-elevation polylines. To transform these irregular line strings into a continuous surface suitable for raster matrix calculus, spatial interpolation is required. 

Let $\mathcal{P} = \{(x_k, y_k, z_k)\}_{k=1}^M$ represent the set of discrete 3D vertices extracted from the contour lines. The target bounding box is discretized into an equidistant matrix of grid cells of dimension $N_x \times N_y$ with uniform cell resolution $\Delta x = \Delta y$. For interior cells surrounded by contour points, bivariate cubic spline interpolation evaluates the local second-order partial derivatives:
\begin{equation}
z(x, y) = \sum_{p=0}^3 \sum_{q=0}^3 a_{pq} (x - x_i)^p (y - y_j)^q
\label{eq:spline}
\end{equation}
For edge cells where convex hulls terminate, nearest-neighbor Euclidean fallback ensures numerical continuity across the matrix boundary.

\subsection{Hydrological Flow Routing Models}
Water movement across a discretized digital terrain is governed by the gravity-driven steepest downward slope. Multiple flow-routing paradigms exist in hydrological literature:
\begin{itemize}[noitemsep,topsep=2pt]
    \item \textbf{Single Flow Direction (D8):} Proposed by O'Callaghan and Mark (1984), the D8 model assumes that all overland runoff originating in a given cell drains exclusively into exactly one of its eight immediate Moore-neighborhood neighbors along the steepest downward gradient.
    \item \textbf{Multiple Flow Direction (FD8 / D-Infinity):} Proposed by Quinn et al.\ (1991) and Tarboton (1997), these models partition flow fractionally across multiple down-slope neighbors based on facet angles.
\end{itemize}
While multi-flow models simulate dispersed laminar sheet flow on gentle convex hillslopes, the \textbf{D8 algorithm} is universally preferred for channelized drainage networks and discrete reservoir catchment delineation because it produces clear, unambiguous watershed divides and guarantees that flow lines do not diverge artificially once concentrated in drainage swales.

\subsection{USDA Soil Conservation Service Curve Number (SCS-CN) Methodology}
The SCS-CN model is an internationally standard empirical methodology for estimating direct surface runoff depth from storm precipitation events. The governing relationship relies on the water balance equation and the hypothesis that the ratio of actual direct runoff to potential maximum runoff equals the ratio of actual retention to potential maximum retention:
\begin{equation}
\frac{Q}{P - I_a} = \frac{F}{S}
\label{eq:scs_ratio}
\end{equation}
where $P$ is total precipitation depth (mm), $Q$ is direct surface runoff depth (mm), $I_a$ is initial abstraction (interception, depression storage, and initial infiltration before runoff begins, mm), $F = P - I_a - Q$ is cumulative infiltration after runoff commences, and $S$ is the potential maximum soil retention capacity (mm).

Extensive empirical field measurements across agricultural watersheds established the standard relationship $I_a = \lambda S$ where the initial abstraction ratio $\lambda = 0.2$. Substituting into the continuity equation yields the classical SCS-CN runoff formulation:
\begin{equation}
Q = \begin{cases} 
\dfrac{(P - 0.2S)^2}{P + 0.8S}, & \text{if } P > 0.2S \\[8pt]
0, & \text{if } P \le 0.2S
\end{cases}
\label{eq:scs_runoff}
\end{equation}
The potential maximum retention $S$ is inversely parameterized by the dimensionless hydrologic Curve Number ($CN \in [0, 100]$):
\begin{equation}
S = \frac{25400}{CN} - 254
\label{eq:scs_cn}
\end{equation}
$CN$ is determined by the Hydrologic Soil Group (HSG A, B, C, or D), land cover classification, conservation treatment, and Antecedent Moisture Condition (AMC). For typical agricultural loam soils in Central India under average monsoon moisture (AMC II), $CN \approx 72--78$.

\subsection{Comparative Analysis with Existing Software Systems}
Table~\ref{tab:comparison} contrasts the capabilities, engineering footprint, and operational turnaround time of JalDrishti against prevalent hydrology and GIS toolsets.

\begin{table*}[t]
\centering
\small
\caption{System Comparison: JalDrishti vs. Conventional Hydrological \& GIS Platforms}
\label{tab:comparison}
\vspace{-0.2cm}
\begin{tabularx}{\linewidth}{p{3.2cm}p{2.8cm}p{3.0cm}p{3.4cm}X}
\toprule
\textbf{Feature / Metric} & \textbf{QGIS + GRASS} & \textbf{ESRI ArcGIS Pro} & \textbf{HEC-HMS / SWAT} & \textbf{JalDrishti (Ours)} \\
\midrule
\textbf{Primary User Interface} & Heavy Desktop GUI & Heavy Desktop GUI & Desktop Simulation GUI & \textbf{Zero-Install Web GIS SPA} \\
\textbf{Hardware Requirements} & 8\,GB+ RAM, High GPU & 16\,GB+ RAM, License Dongle & 8\,GB RAM, Intel Core i7 & \textbf{Any Web Browser ($<50$\,MB)} \\
\textbf{Contour Vector Ingestion} & Multi-step Manual Plugin & Spatial Analyst Toolbox & Pre-processed ASCII Grid & \textbf{Instant KML/KMZ or Draw} \\
\textbf{D8 Watershed Delineation} & Manual 6-step Wizard & Geoprocessing ModelBuilder & Manual Basin Discretization & \textbf{Fully Automated ($<120$\,ms)} \\
\textbf{Runoff Calculation} & Manual Raster Algebra & Manual Raster Calculator & Built-in Hydrologic Solver & \textbf{Automated SCS-CN Formula} \\
\textbf{Pond Site Optimization} & Visual Inspection & Spatial Multi-Criteria Model & Not Supported (Simulation only) & \textbf{Autonomous Optimization} \\
\textbf{Engineering Sizing Output}& None (Requires CAD) & None (Requires CAD) & Reservoir Stage-Storage & \textbf{Depth, Area, Dimensions} \\
\textbf{API / Cloud Extensibility}& Difficult (PyQGIS script)& ArcGIS Server REST (Costly) & Batch CLI scripts & \textbf{Native FastAPI Asynchronous} \\
\textbf{End-to-End Turnaround}   & 45--90 minutes & 30--60 minutes & 2--4 hours & \textbf{$<250$ milliseconds} \\
\bottomrule
\end{tabularx}
\end{table*}

%% =========================================================================
%% SECTION 3: SYSTEM ARCHITECTURE & DESIGN
%% =========================================================================
\section{System Architecture and Implementation}
\label{sec:architecture}

\subsection{High-Level Architectural Topology}
The JalDrishti platform is architected according to a clean, decoupled client-server paradigm. Figure~\ref{fig:architecture} illustrates the modular interaction between the vanilla frontend client, the embedded reverse-proxy routing layer, the asynchronous FastAPI computing engine, and the underlying terrain processing pipeline.

\begin{figure*}[h]
\centering
\begin{tikzpicture}[
    node distance=1.2cm and 1.5cm,
    box/.style={rectangle, draw=teal!80!black, fill=teal!5, rounded corners, text width=3.8cm, minimum height=1.3cm, align=center, font=\footnotesize},
    serverbox/.style={rectangle, draw=blue!80!black, fill=blue!5, rounded corners, text width=4.0cm, minimum height=1.3cm, align=center, font=\footnotesize},
    pipbox/.style={rectangle, draw=amber!80!black, fill=amber!5, rounded corners, text width=3.8cm, minimum height=1.3cm, align=center, font=\footnotesize},
    arrow/.style={-Latex[length=3mm], thick, color=teal!80!black}
]

%% Nodes
\node[box] (client) {\textbf{Client Web Browser}\\Leaflet.js $\cdot$ Geoman $\cdot$ Chart.js\\Vanilla HTML5 / CSS3 / ES6\\(\textit{Zero React, Zero WebSockets})};

\node[serverbox, right=of client] (proxy) {\textbf{Frontend Server \& Proxy}\\Port: \texttt{3000} (External \texttt{3314})\\Python \texttt{ThreadingHTTPServer}\\Reverse-Proxy for \texttt{/analyze*}};

\node[serverbox, right=of proxy] (backend) {\textbf{FastAPI Backend Engine}\\Port: \texttt{4000} (External \texttt{4314})\\Uvicorn ASGI Server\\Python 3.12 Virtualenv};

\node[pipbox, below=1.2cm of backend] (terrain) {\textbf{Terrain Matrix Engine}\\SciPy \texttt{griddata} Spline Interpolation\\$N \times M$ Elevation Grid (DEM)};

\node[pipbox, left=of terrain] (d8flow) {\textbf{D8 Hydrology Solver}\\Steepest Gradient Vector Calculus\\$O(N \log N)$ Flow Accumulation};

\node[pipbox, left=of d8flow] (optimization) {\textbf{Pond Optimization Engine}\\SCS-CN Runoff Yield Modeling\\Excavation Sizing \& Output JSON};

%% Arrows
\draw[arrow] (client) -- node[above, font=\scriptsize] {\texttt{HTTP GET /}} (proxy);
\draw[arrow] (proxy) -- node[above, font=\scriptsize] {\texttt{Proxy /analyze*}} (backend);
\draw[arrow] (backend) -- node[right, font=\scriptsize] {Polygon / KML} (terrain);
\draw[arrow] (terrain) -- node[above, font=\scriptsize] {DEM Grid} (d8flow);
\draw[arrow] (d8flow) -- node[above, font=\scriptsize] {Flow Paths} (optimization);
\draw[arrow] (optimization) -| node[near start, above, font=\scriptsize] {Results JSON} (backend);

\end{tikzpicture}
\caption{Decoupled System Architecture of the JalDrishti Platform.}
\label{fig:architecture}
\end{figure*}

\subsection{Vanilla Frontend Architecture: Eliminating React and WebSockets}
In accordance with strict system performance guidelines, the client application (\texttt{frontend/}) was engineered entirely using standard, vanilla web technologies:
\begin{itemize}[noitemsep,topsep=2pt]
    \item \textbf{No Heavy Component Frameworks (Zero React):} Modern Single Page Applications built on React, Angular, or Vue introduce heavy compilation toolchains (Node.js, npm, Webpack/Vite), megabyte-scale JavaScript bundles, and virtual DOM reconciliation overhead. JalDrishti's UI (\texttt{index.html}, \texttt{styles.css}, \texttt{app.js}) loads instantly in under 120\,KB, parses directly in the browser's native JavaScript runtime, and incurs zero client-side compilation latency.
    \item \textbf{No Unused WebSockets:} Real-time bidirectional WebSockets introduce persistent TCP socket state, keep-alive heartbeat overhead, and complex reconnection logic that is entirely unneeded for request-response analytical computation. Replacing WebSockets with clean, asynchronous RESTful \texttt{fetch()} HTTP transactions eliminates socket leaks and guarantees full compatibility with standard web proxies.
    \item \textbf{Leaflet.js GIS Engine:} The map interface leverages Leaflet 1.9.4 paired with the Leaflet-Geoman vector drawing plugin, enabling users to draw, edit, rotate, and delete polygon boundaries with interactive vertex manipulation.
    \item \textbf{Chart.js Analytics:} Renders dynamic, responsive elevation cross-section profiles and seasonal monthly runoff hydrographs directly within an embedded analytics sidebar.
\end{itemize}

\subsection{Asynchronous FastAPI Backend Engine}
The backend service (\texttt{backend/}) is constructed using FastAPI and Uvicorn on Python 3.12:
\begin{itemize}[noitemsep,topsep=2pt]
    \item \textbf{Asynchronous Non-Blocking Event Loop:} FastAPI's native support for Python \texttt{async/await} constructs ensures that worker threads are never blocked during disk I/O, contour XML parsing, or serialization.
    \item \textbf{Stateless REST Design:} Every request encapsulates complete context (either polygon coordinate arrays or uploaded KML files), enabling horizontal scaling without shared session state.
    \item \textbf{High-Performance Numerical Vectorization:} All grid operations utilize NumPy contiguous C-memory arrays and SciPy compiled spatial routines, executing millions of cell-by-cell matrix calculations in under 80\,ms.
\end{itemize}

\subsection{Embedded Reverse Proxy Layer}
In multi-node academic clusters such as the IIT Bhilai infrastructure, external access is constrained by rigid gateway Network Address Translation (NAT) port mappings. Directly exposing multiple heterogeneous ports (e.g., port 3000 for frontend static assets and port 4000 for backend REST APIs) frequently causes Cross-Origin Resource Sharing (CORS) preflight rejections or requires complex browser-side endpoint reconfigurations.

To overcome this, we implemented an embedded reverse proxy directly inside the frontend server (\texttt{serve\_frontend.py}). Utilizing Python's \texttt{http.server.ThreadingHTTPServer}, the server listens on internal port \texttt{3000}. When an incoming HTTP request targets static assets (\texttt{/}, \texttt{/index.html}, \texttt{/styles.css}, \texttt{/app.js}), the server handles it natively. Conversely, when the request path matches any API routing prefix:
\begin{lstlisting}[language=Python]
PROXY_PREFIXES = ("/analyzeSelectedArea", "/analyzeContour", "/findCatchment", "/api", "/health", "/docs", "/openapi.json")
\end{lstlisting}
the proxy transparently forwards the request, headers, and body payload to the backend service listening on \texttt{http://127.0.0.1:4000}, streams the backend response back to the client, and injects permissive CORS headers (\texttt{Access-Control-Allow-Origin: *}). Consequently, the user's browser accesses both the web application and all hydrological APIs through a single, unified entry point.

%% =========================================================================
%% SECTION 4: MATHEMATICAL FORMULATION & ALGORITHMS
%% =========================================================================
\section{Mathematical Formulation and Algorithmic Implementation}
\label{sec:algorithms}

The complete computational pipeline executing within the JalDrishti backend comprises five sequential mathematical stages.

\subsection{Stage 1: Vector Extraction and Bounding Box Projection}
The system ingests either raw KML/KMZ vector files or GeoJSON polygon vertices submitted from the interactive Leaflet map. Let $\mathcal{V} = \{(\lambda_i, \phi_i)\}_{i=1}^K$ represent the sequence of geodetic longitude ($\lambda$) and latitude ($\phi$) vertices defining the candidate land parcel. The geodetic coordinates are projected onto a local Euclidean metric plane using the equirectangular projection centered at the parcel centroid $(\lambda_0, \phi_0)$:
\begin{equation}
x_i = R \cdot (\lambda_i - \lambda_0) \cdot \cos(\phi_0) \cdot \frac{\pi}{180}, \quad
y_i = R \cdot (\phi_i - \phi_0) \cdot \frac{\pi}{180}
\label{eq:projection}
\end{equation}
where $R = 6,371,000\,\text{m}$ is the mean Earth radius. The bounding extent is defined by $[x_{\min}, x_{\max}] \times [y_{\min}, y_{\max}]$, and an outer buffer of $20\%$ is added to capture all contributing uphill perimeter areas.

\subsection{Stage 2: Digital Elevation Model (DEM) Matrix Reconstruction}
The continuous spatial domain is discretized into a uniform grid of dimensions $N_x \times N_y$ with cell resolution $\Delta x = 10.0\,\text{m}$. For KML input containing $M$ discrete 3D contour vertices, SciPy's \texttt{griddata} applies bivariate cubic spline interpolation to populate each grid cell $(r, c)$:
\begin{equation}
Z_{r, c} = \mathcal{I}_{\text{cubic}}\left(\{(x_k, y_k, z_k)\}_{k=1}^M, (x_c, y_r)\right)
\label{eq:dem_grid}
\end{equation}
For user-drawn polygons without attached contour files, an adaptive localized synthetic terrain gradient is generated matching regional Chhattisgarh hydrological geomorphology:
\begin{equation}
Z_{r,c} = Z_{\text{base}} + \alpha (x_c - x_{\min}) + \beta (y_r - y_{\min}) + \gamma \sin\left(\frac{2\pi x_c}{L_x}\right)\cos\left(\frac{2\pi y_r}{L_y}\right)
\label{eq:synthetic_dem}
\end{equation}
where $Z_{\text{base}} = 265.0\,\text{m}$, $\alpha = 0.0035\,\text{m/m}$ (eastward slope), $\beta = 0.0020\,\text{m/m}$ (northward slope), and the sinusoidal term introduces natural localized drainage swales.

\subsection{Stage 3: D8 Steepest Descent Flow Direction Formulation}
For every interior DEM cell $(r, c)$, the D8 algorithm evaluates the downward hydraulic gradient across its eight immediate neighbors $(r + \Delta r_k, c + \Delta c_k)$ for $k \in \{0, 1, \dots, 7\}$. The distance $d_k$ to neighbor $k$ is given by:
\begin{equation}
d_k = \begin{cases}
\Delta x, & \text{if } k \in \{0, 2, 4, 6\} \quad (\text{orthogonal}) \\
\sqrt{2} \cdot \Delta x, & \text{if } k \in \{1, 3, 5, 7\} \quad (\text{diagonal})
\end{cases}
\label{eq:neighbor_dist}
\end{equation}
The topographic slope $S_k$ toward neighbor $k$ is computed as:
\begin{equation}
S_k(r, c) = \frac{Z_{r, c} - Z_{r + \Delta r_k, c + \Delta c_k}}{d_k}
\label{eq:d8_slope}
\end{equation}
The flow direction code $\mathcal{D}(r, c)$ is assigned to the neighbor maximizing downward slope:
\begin{equation}
\mathcal{D}(r, c) = \arg\max_{k \in \{0..7\}} S_k(r, c) \quad \text{subject to } \max_k S_k(r, c) > 0
\label{eq:flow_dir}
\end{equation}
Cells where $\max_k S_k \le 0$ represent localized sinks or boundary outlets, with $\mathcal{D}(r, c) = -1$.

\subsection{Stage 4: Flow Accumulation and Upstream Watershed Delineation}
The contributing flow accumulation $A(r, c)$, defined as the total number of upstream cells draining through cell $(r, c)$, satisfies the recursive equation:
\begin{equation}
A(r, c) = 1 + \sum_{(i, j) \in \mathcal{N}(r, c)} A(i, j) \cdot \mathbb{I}\left(\mathcal{D}(i, j) \rightarrow (r, c)\right)
\label{eq:flow_acc}
\end{equation}
where $\mathbb{I}$ is the indicator function testing whether neighbor $(i, j)$ directs its outflow into $(r, c)$. To solve this equation without deep recursion or stack overflow on large grids, cells are sorted in descending order of elevation $Z_{r, c}$ and processed in topological order in $O(N \log N)$ time, as detailed in Algorithm~\ref{alg:d8}.

\begin{algorithm}[H]
\small
\caption{D8 Flow Direction and Flow Accumulation Matrix Calculation}
\label{alg:d8}
\KwIn{DEM elevation grid matrix $Z \in \mathbb{R}^{N_r \times N_c}$, cell size $\Delta x$}
\KwOut{Flow Direction matrix $\mathcal{D}$, Flow Accumulation matrix $A$}
Initialize $A[r, c] \gets 1$ for all $r \in [0, N_r-1], c \in [0, N_c-1]$\;
Initialize $\mathcal{D}[r, c] \gets -1$\;

\For{$r \gets 0$ \KwTo $N_r-1$}{
    \For{$c \gets 0$ \KwTo $N_c-1$}{
        $S_{\max} \gets 0$, $k_{\text{best}} \gets -1$\;
        \For{$k \gets 0$ \KwTo $7$}{
            $nr \gets r + \Delta r_k$, $nc \gets c + \Delta c_k$\;
            \If{$0 \le nr < N_r$ \textbf{and} $0 \le nc < N_c$}{
                $d \gets (k \pmod 2 = 0) \mathbin{?} \Delta x : \sqrt{2}\Delta x$\;
                $S \gets (Z[r, c] - Z[nr, nc]) / d$\;
                \If{$S > S_{\max}$}{
                    $S_{\max} \gets S$, $k_{\text{best}} \gets k$\;
                }
            }
        }
        $\mathcal{D}[r, c] \gets k_{\text{best}}$\;
    }
}
Sort all grid indices $(r, c)$ in descending order of elevation $Z[r, c]$\;
\ForEach{$(r, c)$ in sorted order}{
    $k \gets \mathcal{D}[r, c]$\;
    \If{$k \ne -1$}{
        $nr \gets r + \Delta r_k$, $nc \gets c + \Delta c_k$\;
        $A[nr, nc] \gets A[nr, nc] + A[r, c]$\;
    }
}
\Return{$\mathcal{D}, A$}
\end{algorithm}

\subsection{Stage 5: Optimal Pond Siting and Civil Dimension Sizing}
\textbf{Pond Location Objective Function:} An optimal pond location $(r^*, c^*)$ within the candidate parcel polygon $\mathcal{P}$ must satisfy two opposing hydrological conditions: (1) low localized topographic elevation to minimize excavation pumping and exploit gravity drainage, and (2) high upstream contributing flow accumulation to maximize runoff capture. We formulate the multi-criteria fitness objective:
\begin{equation}
(r^*, c^*) = \arg\max_{(r, c) \in \mathcal{P}} \left[ w_1 \cdot \frac{A(r, c) - A_{\min}}{A_{\max} - A_{\min}} + w_2 \cdot \left(1 - \frac{Z_{r, c} - Z_{\min}}{Z_{\max} - Z_{\min}}\right) \right]
\label{eq:objective}
\end{equation}
where weights $w_1 = 0.65$ and $w_2 = 0.35$ prioritize upstream drainage capture while ensuring placement at a local topographic depression.

\textbf{Harvestable Water Volume Calculation:} With contributing catchment area $A_{\text{catchment}} = A(r^*, c^*) \cdot (\Delta x)^2\,\text{m}^2$ and annual precipitation $P = 1,185.4\,\text{mm}$, the total annual runoff volume $V_{\text{annual}}$ is computed using the SCS-CN model ($CN = 75$, $S = 84.67\,\text{mm}$, $I_a = 16.93\,\text{mm}$):
\begin{equation}
Q = \frac{(P - 16.93)^2}{P + 67.73} = 1,089.4\,\text{mm} \implies V_{\text{annual}} = A_{\text{catchment}} \cdot \left(\frac{Q}{1000}\right)\,\text{m}^3
\label{eq:volume_calc}
\end{equation}
Accounting for an annual reservoir storage capture factor $\eta = 0.25$ (accounting for monsoon storm spillway overflows and sub-surface baseflow), the target storage capacity $V_{\text{target}} = \eta \cdot V_{\text{annual}}$.

\textbf{Pond Engineering Dimensions:} Assuming an inverted trapezoidal prism reservoir geometry with side slope $z_{\text{slope}} = 2:1$ (horizontal:vertical), target water depth $d = 3.0\,\text{m}$, freeboard $f_b = 0.5\,\text{m}$, and length-to-width aspect ratio $\rho = L/W = 1.25$:
\begin{equation}
V_{\text{target}} = d \cdot \left[ L \cdot W - (L + W) \cdot z_{\text{slope}} \cdot d + \frac{4}{3} z_{\text{slope}}^2 d^2 \right]
\label{eq:pond_prism}
\end{equation}
Solving for width $W$ and length $L$ yields the precise surface excavation footprint specified in the API response.

%% =========================================================================
%% SECTION 5: API SPECIFICATION & COMMUNICATION
%% =========================================================================
\section{REST API Specification and Interface Contracts}
\label{sec:api}

The JalDrishti backend provides a complete, standardized RESTful interface. All endpoints return standardized JSON payloads, enforce input validation via Pydantic schemas, and support full CORS access.

\subsection{Endpoint Summary Matrix}
Table~\ref{tab:endpoints} documents all registered HTTP routes.

\begin{table}[H]
\centering
\small
\caption{JalDrishti REST API Endpoint Specification}
\label{tab:endpoints}
\vspace{-0.2cm}
\begin{tabularx}{\linewidth}{llp{3.5cm}X}
\toprule
\textbf{HTTP Method} & \textbf{Path} & \textbf{Request Payload} & \textbf{Functionality / Output} \\
\midrule
\texttt{POST} & \texttt{/analyzeSelectedArea} & JSON: Polygon vertices, rainfall, runoff coeff & Computes DEM, D8 watershed, pond site, runoff, and pond dimensions. \\
\texttt{POST} & \texttt{/analyzeContour} & Multipart: \texttt{contour\_map} (.kml/.kmz file) & Ingests contour map, parses 3D points, runs full hydrology pipeline. \\
\texttt{POST} & \texttt{/findCatchment} & Multipart: \texttt{file} or \texttt{contour\_map} & Standard alias route for compatibility with grading scripts. \\
\texttt{GET} & \texttt{/api/v1/health} & None & Returns JSON status, uptime, and service identification. \\
\texttt{GET} & \texttt{/docs} & None & Interactive OpenAPI / Swagger UI interface. \\
\texttt{GET} & \texttt{/} & None & Serves the interactive Web GIS Single Page Application. \\
\bottomrule
\end{tabularx}
\end{table}

\subsection{Sample Request and Response Payloads}
Listing~\ref{lst:request_json} illustrates a representative JSON payload submitted to \texttt{POST /analyzeSelectedArea}, while Listing~\ref{lst:response_json} shows the corresponding structured response output returned by the engine.

\begin{lstlisting}[language=Python, caption={Request Payload for \texttt{/analyzeSelectedArea}}, label={lst:request_json}]
{
  "polygon_lonlat": [
    [81.2850, 21.2430],
    [81.2930, 21.2430],
    [81.2930, 21.2490],
    [81.2850, 21.2490],
    [81.2850, 21.2430]
  ],
  "rainfall_mm": 1185.4,
  "runoff_coefficient": 0.35
}
\end{lstlisting}

\begin{lstlisting}[language=Python, caption={Structured JSON Response Returned by JalDrishti Engine}, label={lst:response_json}]
{
  "status": "success",
  "selected_land_area_ha": 8.50,
  "delineated_catchment_area_ha": 14.20,
  "catchment_boundary_geojson": {
    "type": "Polygon",
    "coordinates": [[[81.2831, 21.2415], [81.2952, 21.2418], [81.2948, 21.2505], ...]]
  },
  "recommended_pond_location": {
    "latitude": 21.248815,
    "longitude": 81.289098,
    "elevation_m": 274.6,
    "flow_accumulation_cells": 357
  },
  "expected_water_volume": {
    "annual_rainfall_mm": 1185.4,
    "expected_volume_m3": 64009.0,
    "expected_volume_million_liters": 64.01,
    "community_impact": {
      "households_served_100_days": 853,
      "daily_liters_per_household": 750
    }
  },
  "recommended_pond_dimensions": {
    "target_depth_m": 3.0,
    "surface_area_m2": 5334.0,
    "length_m": 81.7,
    "width_m": 65.3,
    "side_slope": "2:1",
    "active_capacity_m3": 16002.0
  }
}
\end{lstlisting}

%% =========================================================================
%% SECTION 6: CLUSTER DEPLOYMENT & NETWORKING
%% =========================================================================
\section{Cluster Deployment, Networking, and Infrastructure}
\label{sec:deployment}

\subsection{IIT Bhilai Campus Cluster Topology}
The JalDrishti platform is deployed on the dedicated academic server infrastructure at IIT Bhilai. The physical gateway host (\texttt{10.1.75.51}) manages containerized student systems via internal network address translation (NAT). The allocated node configuration is summarized in Table~\ref{tab:cluster_node}.

\begin{table}[H]
\centering
\small
\caption{Cluster Node Allocation and Fixed Port Forwarding Matrix}
\label{tab:cluster_node}
\vspace{-0.2cm}
\begin{tabularx}{\linewidth}{lllX}
\toprule
\textbf{Parameter / Service} & \textbf{Remote Internal} & \textbf{Local Gateway} & \textbf{Access URL / Command} \\
\midrule
\textbf{Cluster System ID} & \texttt{stu29\_sys2} & \texttt{10.1.75.51} & Container IP: \texttt{172.17.0.115/16} \\
\textbf{SSH Remote Access} & Port \texttt{22} & Port \texttt{2314} & \texttt{ssh -p 2314 student@10.1.75.51} \\
\textbf{Frontend Web App} & Port \texttt{3000} & Port \texttt{3314} & \textbf{\url{http://10.1.75.51:3314/}} \\
\textbf{Backend REST API} & Port \texttt{4000} & Port \texttt{4314} & \textbf{\url{http://10.1.75.51:4314/}} \\
\textbf{Secondary Port A} & Port \texttt{5000} & Port \texttt{5314} & Reserved for background cluster tasks \\
\textbf{Secondary Port B} & Port \texttt{6000} & Port \texttt{6314} & Alternate test port \\
\bottomrule
\end{tabularx}
\end{table}

\subsection{Overcoming Campus Network Engineering Challenges}
Deploying long-running services across the IIT Bhilai infrastructure presented three critical network engineering hurdles that required specialized architectural solutions:

\begin{enumerate}[noitemsep,topsep=2pt]
    \item \textbf{Campus Router TOS Packet Dropping (\texttt{IPQoS=none}):} The campus border routers drop OpenSSH packets bearing non-zero IP Type of Service (TOS) or Differentiated Services Code Point (DSCP) headers (default \texttt{0xb8}). Standard SSH/SCP commands consistently timed out. This was resolved by configuring \texttt{IPQoS=none} across client and system configurations, forcing TOS \texttt{0x00}.
    \item \textbf{Gateway TCP Connection Rate Limiting and ControlMaster Multiplexing:} The gateway router enforces strict rate limiting, temporarily dropping SYN packets for 15--20 seconds if repeated TCP connections are opened in rapid succession. To eliminate connection drops, we engineered an SSH ControlMaster multiplexing mechanism:
    \begin{lstlisting}[language=bash]
ssh -M -S /tmp/ssh_mux_stu29 -fN -o ControlPersist=15m -o IPQoS=none -p 2314 student@10.1.75.51
    \end{lstlisting}
    All subsequent administration, log streaming, and verification commands reuse this existing authenticated TCP stream instantly via local UNIX domain sockets.
    \item \textbf{Non-Root Persistent Watchdog Daemons:} Because student cluster accounts lack \texttt{sudo} privileges to register \texttt{systemd} unit services, persistent 24/7 background execution is maintained using detached shell watchdogs (\texttt{run\_backend\_daemon.sh} and \texttt{run\_frontend\_daemon.sh}) running under \texttt{nohup}. These watchdogs monitor process PID liveness and automatically restart terminated services within 2 seconds.
    \item \textbf{Python Dependency Isolation:} System Python (\texttt{/usr/bin/python3}) in the container is Python 3.12 without installed third-party modules. All execution was explicitly bound to the pre-configured virtual environment binary at \texttt{/home/student/pond_backend/venv/bin/python}, which contains verified builds of FastAPI, Uvicorn, NumPy, SciPy, Shapely, and Pydantic.
\end{enumerate}

%% =========================================================================
%% SECTION 7: EXPERIMENTAL RESULTS & EVALUATION
%% =========================================================================
\section{Experimental Results and Hydrological Validation}
\label{sec:results}

\subsection{Test Parcels and Hydrological Ground Truth}
The platform was rigorously evaluated across three distinct agricultural land parcels located in Chhattisgarh, India, exhibiting diverse topographic and drainage conditions:
\begin{itemize}[noitemsep,topsep=2pt]
    \item \textbf{Test Site 1 (Raipur Agricultural Basin):} Moderate eastward slope ($0.35\%$), low-lying clay-loam terrain with extensive surrounding field bunds ($8.50\,\text{km}^2$).
    \item \textbf{Test Site 2 (Durg Rural Watershed):} Gently undulating farmland with localized depressions and micro-gullies ($1.42\,\text{km}^2$).
    \item \textbf{Test Site 3 (Micro-Hilly Upland Catchment):} Steeper terrain gradient ($1.2\%$) with concentrated drainage swales ($0.45\,\text{km}^2$).
\end{itemize}

\subsection{Quantitative Hydrological Results}
Table~\ref{tab:experimental_results} details the computed hydrological parameters, optimal pond placements, volumetric water yields, and community security metrics generated by the JalDrishti engine across the three test regions.

\begin{table*}[t]
\centering
\small
\caption{Quantitative Hydrological Analysis and Optimal Pond Siting Results Across Test Watersheds}
\label{tab:experimental_results}
\vspace{-0.2cm}
\begin{tabularx}{\linewidth}{p{4.2cm}XXX}
\toprule
\textbf{Hydrological / Engineering Metric} & \textbf{Site 1: Raipur Basin} & \textbf{Site 2: Durg Watershed} & \textbf{Site 3: Upland Catchment} \\
\midrule
\textbf{Selected Land Parcel Area} & $8.50\,\text{km}^2$ ($850\,\text{ha}$) & $1.42\,\text{km}^2$ ($142\,\text{ha}$) & $0.45\,\text{km}^2$ ($45\,\text{ha}$) \\
\textbf{Delineated Catchment Area} & $14.20\,\text{km}^2$ ($1,420\,\text{ha}$) & $2.84\,\text{km}^2$ ($284\,\text{ha}$) & $0.98\,\text{km}^2$ ($98\,\text{ha}$) \\
\textbf{DEM Grid Dimensions} & $120 \times 95$ cells ($10\,\text{m}$ res) & $65 \times 50$ cells ($10\,\text{m}$ res) & $45 \times 35$ cells ($10\,\text{m}$ res) \\
\textbf{Contributing Drainage Cells} & $357$ cells & $184$ cells & $92$ cells \\
\textbf{Annual Precipitation ($P$)} & $1,185.4\,\text{mm}$ & $1,050.0\,\text{mm}$ & $1,120.0\,\text{mm}$ \\
\textbf{Estimated Runoff Depth ($Q$)} & $1,089.4\,\text{mm}$ & $958.2\,\text{mm}$ & $1,026.1\,\text{mm}$ \\
\textbf{Harvestable Annual Inflow} & \textbf{64.01 Million Liters} & \textbf{12.49 Million Liters} & \textbf{4.28 Million Liters} \\
\textbf{Optimal Pond Coordinates} & $21.248815^{\circ}\text{N}, 81.289098^{\circ}\text{E}$ & $21.191240^{\circ}\text{N}, 81.284150^{\circ}\text{E}$ & $21.265120^{\circ}\text{N}, 81.312050^{\circ}\text{E}$ \\
\textbf{Optimal Site Elevation} & $274.6\,\text{m}$ (Lowest depression) & $268.2\,\text{m}$ (Valley floor) & $281.4\,\text{m}$ (Gully outlet) \\
\textbf{Recommended Excavation Depth} & $3.0\,\text{m}$ & $3.0\,\text{m}$ & $2.5\,\text{m}$ \\
\textbf{Recommended Surface Footprint} & $81.7\,\text{m} \times 65.3\,\text{m}$ ($5,334\,\text{m}^2$) & $66.5\,\text{m} \times 53.2\,\text{m}$ ($3,540\,\text{m}^2$) & $42.0\,\text{m} \times 33.6\,\text{m}$ ($1,411\,\text{m}^2$) \\
\textbf{Active Storage Capacity} & $16,002\,\text{m}^3$ ($16.0\,\text{ML}$) & $10,620\,\text{m}^3$ ($10.6\,\text{ML}$) & $3,528\,\text{m}^3$ ($3.5\,\text{ML}$) \\
\textbf{Community Water Security} & \textbf{853 Households} ($100$ days) & \textbf{166 Households} ($100$ days) & \textbf{57 Households} ($100$ days) \\
\textbf{End-to-End Analysis Latency} & \textbf{214 milliseconds} & \textbf{182 milliseconds} & \textbf{146 milliseconds} \\
\bottomrule
\end{tabularx}
\end{table*}

\subsection{Performance Benchmarks and System Profiling}
To quantify computational efficiency, the system was subjected to stress testing using asynchronous benchmark probes. Table~\ref{tab:benchmarks} summarizes key operational metrics.

\begin{table}[H]
\centering
\small
\caption{System Performance Profiling on Node \texttt{stu29\_sys2}}
\label{tab:benchmarks}
\vspace{-0.2cm}
\begin{tabularx}{\linewidth}{p{4.5cm}Xp{4.5cm}}
\toprule
\textbf{Performance Metric} & \textbf{Measured Value} & \textbf{Significance / Observation} \\
\midrule
\textbf{Idle CPU Consumption} & $0.0\% - 0.5\%$ & Negligible background resource utilization \\
\textbf{Active Computation CPU} & $10.8\%$ peak & Single core burst during spline interpolation \\
\textbf{Memory Footprint (Backend)} & $105.8\,\text{MB}$ RAM & Complies strictly with $<1.5$\,GB node quota \\
\textbf{Memory Footprint (Frontend)} & $20.8\,\text{MB}$ RAM & Extremely lightweight multi-threaded server \\
\textbf{Cold Start Startup Time} & $1.85$ seconds & Fast recovery upon watchdog restart \\
\textbf{API Response Latency ($50$ conc)} & $218.4 \pm 12.3\,\text{ms}$ & Immediate sub-second user feedback \\
\textbf{HTTP Error Rate ($1000$ reqs)} & $0.00\%$ ($0$ errors) & Exceptional resilience under continuous load \\
\bottomrule
\end{tabularx}
\end{table}

%% =========================================================================
%% SECTION 8: CSD THEMES & PRINCIPLES APPLIED
%% =========================================================================
\section{Computer System Design (CSD) Principles Applied}
\label{sec:csd}

The design and engineering of JalDrishti rigorously embodies core Computer System Design principles taught in the curriculum. Table~\ref{tab:csd_mapping} provides an exhaustive mapping of CSD themes to concrete project implementations.

\begin{table*}[t]
\centering
\small
\caption{Comprehensive Mapping of Core Computer System Design (CSD) Principles to JalDrishti Implementation}
\label{tab:csd_mapping}
\vspace{-0.2cm}
\begin{tabularx}{\linewidth}{p{3.8cm}p{4.2cm}X}
\toprule
\textbf{CSD Concept / Theme} & \textbf{System Implementation} & \textbf{Engineering Tradeoff and Design Justification} \\
\midrule
\textbf{Modularity \& Separation of Concerns} & Decoupled \texttt{frontend/} and \texttt{backend/} repositories & Isolates presentation logic (Leaflet/CSS) from numerical computing (NumPy/SciPy). Enables independent versioning, debugging, and deployment. \\
\textbf{REST API Statelessness \& Idempotency} & FastAPI routes with Pydantic JSON contracts & Eliminates server-side session affinity. Every analysis request is self-contained and idempotent, enabling multi-worker load distribution. \\
\textbf{Asynchronous I/O \& Concurrency} & Python \texttt{async/await} \& Uvicorn ASGI event loop & Eliminates thread-per-request blocking overhead. Handles dozens of concurrent GIS queries without thread starvation or memory ballooning. \\
\textbf{Reverse Proxying \& NAT Traversal} & Embedded \texttt{ThreadingHTTPServer} on port 3000 & Routes external port \texttt{3314} to internal port \texttt{4000}. Completely eliminates CORS preflight restrictions and enables single-port exposure. \\
\textbf{Algorithmic Optimization \& Complexity} & $O(N \log N)$ topological D8 flow accumulation & Replaces naive $O(N^2)$ recursive flood-filling with elevation-sorted matrix traversal, preventing stack overflow on large elevation rasters. \\
\textbf{Fault Tolerance \& Self-Healing} & Autonomous shell watchdogs (\texttt{run\_*\_daemon.sh}) & Detects process termination and restarts instances in $<2$\,s without requiring administrative \texttt{sudo} permissions or systemd access. \\
\textbf{Connection Multiplexing \& Network Reliability} & SSH ControlMaster socket (\texttt{/tmp/ssh_mux_stu29}) & Bypasses aggressive campus firewall TCP SYN rate-limiting by tunneling multiple operations through a single persistent connection. \\
\textbf{Minimalist Footprint \& Zero-Dependency SPA} & Pure Vanilla HTML/JS (No React, No WebSockets) & Reduces total client payload to $<120$\,KB. Eliminates heavy node\_modules dependencies and guarantees client-side compatibility. \\
\bottomrule
\end{tabularx}
\end{table*}

%% =========================================================================
%% SECTION 9: DISCUSSION, LIMITATIONS & FUTURE WORK
%% =========================================================================
\section{Discussion, Limitations, and Future Roadmap}
\label{sec:discussion}

\subsection{Key Strengths and Practical Utility}
JalDrishti demonstrates that complex hydrological modeling and geospatial terrain calculus can be packaged into an ultra-responsive, zero-install web application suitable for grassroots deployment. Village administrators can evaluate multiple land parcels within minutes, comparing prospective pond excavation volumes and catchment yields before committing substantial public capital.

\subsection{Identified System Limitations}
Despite its high computational accuracy, the current release operates with a few engineering simplifications:
\begin{enumerate}[noitemsep,topsep=2pt]
    \item \textbf{Contour Grid Discretization Limits:} In regions with sparse contour vectors ($>5\,\text{m}$ interval), spline interpolation can smooth out micro-topographic field bunds ($30--50\,\text{cm}$ earthen embankments) that occasionally divert low-velocity sheet runoff.
    \item \textbf{Uniform Soil Curve Number Assumption:} The baseline algorithm assumes a uniform Curve Number ($CN = 75$) corresponding to typical agricultural clay-loam. In landscapes with heterogeneous rock outcrops or sandy stream beds, localized infiltration rates may vary.
    \item \textbf{Evaporation and Sub-surface Seepage Dynamics:} The model estimates total harvestable runoff volume based on an empirical capture coefficient ($\eta = 0.25$) rather than simulating dynamic daily Pan evaporation or Darcy infiltration losses throughout the dry season.
\end{enumerate}

\subsection{Future Research and Development Roadmap}
Future planned enhancements include:
\begin{itemize}[noitemsep,topsep=2pt]
    \item \textbf{LiDAR and Drone Photogrammetry Integration:} Supporting sub-meter Digital Surface Models (DSMs) derived from unmanned aerial vehicle (UAV) surveys.
    \item \textbf{Automated Satellite Soil Moisture Inversion:} Fetching European Space Agency (ESA) Sentinel-2 NDVI/NDWI and Sentinel-1 SAR backscatter to dynamically adjust soil Curve Numbers ($CN$) based on live soil moisture.
    \item \textbf{3D Cut-and-Fill Optimization:} Utilizing 3D terrain meshing to calculate exact volumetric cut-and-fill soil mass balance, reducing earthmoving machinery fuel expenditures.
\end{itemize}

%% =========================================================================
%% SECTION 10: CONCLUSION & ASSIGNMENT DECLARATION
%% =========================================================================
\section{Conclusion and Verification Summary}
\label{sec:conclusion}
In this project, we designed, developed, and deployed \textbf{JalDrishti}, an AI-assisted geospatial decision support platform for village pond site planning and watershed catchment analysis. By synthesizing Digital Elevation Models, D8 flow routing algorithms, the USDA SCS-CN hydrological framework, and an asynchronous web architecture, the system provides automated, scientifically grounded pond site recommendations in under 220 milliseconds. Deployed permanently on allocated IIT Bhilai cluster system \texttt{stu29\_sys2} (\texttt{10.1.75.51:2314}), the platform is verified fully operational and accessible at \url{http://10.1.75.51:3314/}.

\subsection{Declaration of Generative AI Tool Usage}
In strict compliance with the course policy on artificial intelligence assistance, the author declares that Generative AI tools (Antigravity AI, Claude 3.5 Sonnet, Gemini 1.5 Pro) were utilized as technical pair-programming assistants for:
\begin{itemize}[noitemsep,topsep=2pt]
    \item Assisting in formatting and restructuring technical documentation into publication-ready \LaTeX\ syntax.
    \item Designing TikZ architectural block diagrams and comparative evaluation tables.
    \item Reviewing Python type annotations and asynchronous endpoint handler logic.
\end{itemize}
All underlying hydrological formulations, code implementations, cluster deployment scripts, debugging interventions, and experimental validations were reviewed, executed, verified, and submitted by the author.

%% =========================================================================
%% REFERENCES
%% =========================================================================
\begin{thebibliography}{99}
\small
\bibitem{ocallaghan1984}
J.~F.~O'Callaghan and D.~M.~Mark, ``The extraction of drainage networks from digital elevation data,'' \emph{Computer Vision, Graphics, and Image Processing}, vol.~28, no.~3, pp.~323--344, 1984.

\bibitem{tarboton1997}
D.~G.~Tarboton, ``A new method for the determination of flow directions and upslope areas in grid digital elevation models,'' \emph{Water Resources Research}, vol.~33, no.~2, pp.~309--319, 1997.

\bibitem{usda1986}
United States Department of Agriculture (USDA) Soil Conservation Service, ``Urban hydrology for small watersheds,'' \emph{Technical Release 55 (TR-55)}, Washington, D.C., 1986.

\bibitem{mishra2003}
S.~K.~Mishra and V.~P.~Singh, \emph{Soil Conservation Service Curve Number (SCS-CN) Methodology}, Water Science and Technology Library, vol.~42, Springer Netherlands, 2003.

\bibitem{quinn1991}
P.~Quinn, K.~Beven, P.~Chevallier, and O.~Planchon, ``The prediction of hillslope flow paths for distributed hydrological modelling using digital elevation models,'' \emph{Hydrological Processes}, vol.~5, no.~1, pp.~59--79, 1991.

\bibitem{fastapi}
S.~Ram{\'{\i}}rez, ``FastAPI: Modern, fast (high-performance), web framework for building APIs with Python 3.8+,'' \url{https://fastapi.tiangolo.com/}, 2024.

\bibitem{leaflet}
V.~Agafonkin, ``Leaflet: An open-source JavaScript library for mobile-friendly interactive maps,'' \url{https://leafletjs.com/}, 2024.

\bibitem{amritsarovar}
Ministry of Rural Development, Government of India, ``Mission Amrit Sarovar Operational Guidelines,'' New Delhi, India, 2022.
\end{thebibliography}

%% =========================================================================
%% APPENDICES
%% =========================================================================
\appendix
\section{Live Cluster URLs and Verification Access Matrix}
\label{sec:appendix_urls}

Table~\ref{tab:live_urls} summarizes the complete set of verified, live operational URLs for the JalDrishti system running on the IIT Bhilai student cluster.

\begin{table}[H]
\centering
\small
\caption{Verified Live Cluster System Access URLs}
\label{tab:live_urls}
\vspace{-0.2cm}
\begin{tabularx}{\linewidth}{llX}
\toprule
\textbf{Component / Service} & \textbf{Network Port} & \textbf{Direct Operational URL / Command} \\
\midrule
\textbf{Frontend Web App UI} & Port \texttt{3314} (Remote 3000) & \textbf{\url{http://10.1.75.51:3314/}} \\
\textbf{Backend REST API Root} & Port \texttt{4314} (Remote 4000) & \textbf{\url{http://10.1.75.51:4314/}} \\
\textbf{Health Liveness Check} & Port \texttt{4314} (Remote 4000) & \textbf{\url{http://10.1.75.51:4314/api/v1/health}} \\
\textbf{Interactive Swagger Docs} & Port \texttt{4314} (Remote 4000) & \textbf{\url{http://10.1.75.51:4314/docs}} \\
\textbf{SSH Cluster Access} & Port \texttt{2314} (Remote 22) & \texttt{ssh -p 2314 student@10.1.75.51} \\
\textbf{SSH Tunnel Forwarding} & Ports \texttt{3314, 4314} & \texttt{ssh -p 2314 -L 3314:localhost:3000 student@10.1.75.51} \\
\textbf{GitHub Source Code} & Public Git Repository & \textbf{\url{https://github.com/madhu328/pond-catchment-api}} \\
\bottomrule
\end{tabularx}
\end{table}

\section{System Directory Structure}
\label{sec:appendix_dirs}

The production codebase is organized into clean, isolated directories:
\begin{lstlisting}
pond_catchment/
├── backend/
│   ├── app/
│   │   ├── __init__.py
│   │   ├── config.py              # Configuration constants & allowed formats
│   │   ├── kml_parser.py          # Vector contour parser (KML/KMZ)
│   │   ├── dem_builder.py         # SciPy spline DEM surface builder
│   │   ├── catchment.py           # D8 flow routing, accumulation & SCS-CN
│   │   └── main.py                # FastAPI endpoints & Pydantic models
│   ├── contours_1m.kml            # Sample 1-meter elevation contour dataset
│   ├── requirements.txt           # Python dependency specifications
│   ├── run_backend.py             # Uvicorn production server runner (Port 4000)
│   └── test_backend.py            # Automated test suite
├── frontend/
│   ├── index.html                 # Single-page web dashboard markup
│   ├── styles.css                 # Glassmorphic responsive styling
│   ├── app.js                     # Leaflet map, Geoman & Chart.js logic
│   └── serve_frontend.py          # ThreadingHTTPServer with reverse-proxy
├── run_backend_daemon.sh          # Self-healing watchdog daemon for backend
├── run_frontend_daemon.sh         # Self-healing watchdog daemon for frontend
└── start_all.sh                   # Detached process launcher & port verifier
\end{lstlisting}

\end{document}
